\documentclass[11pt, a4paper]{article}

\usepackage[UTF8]{ctex}
\usepackage{geometry}
\usepackage{titlesec}
\usepackage{enumitem}
\usepackage{setspace}
\usepackage{xcolor}
\usepackage{hyperref}

\geometry{a4paper, margin=2.2cm}
\setstretch{1.2}

\hypersetup{
	colorlinks=true,
	linkcolor=black,
	urlcolor=blue
}

\title{角色表情图提示词汇总\\（严格参考人物图版）}
\author{黄浩宇}
\date{\today}

\begin{document}
	
	\maketitle
	\newpage
	
	\tableofcontents
	\newpage
	
	\section{统一出图格式规范}
	
	\subsection*{核心原则}
	\begin{itemize}[leftmargin=2em, itemsep=0.2em]
		\item 人物外貌、服装、发型、色彩一律以参考人物图为准，提示词中不重复描述，只做“严格参考”约束。
		\item 一张图 = 同一角色 6 个表情，水平并排一行（1×6 横向排列），等宽等距。
		\item 每个表情为半身像/胸部以上，正面平视，头部居中，构图一致。
		\item 背景：纯色浅灰，无杂物、无文字、无水印。
		\item 风格：写实风格，电影感光影，高清细节，面部特写清晰。
		\item 除表情外，人物长相、服装、发型、光线、机位、比例必须与参考图完全一致，仅表情变化。
		\item 表情差异明显、可辨。
	\end{itemize}
	
	\subsection*{统一前缀（固定加在每条开头）}
	\begin{quote}
		\small
		严格参考所提供的人物图，人物长相、服装、发型、色彩、比例与参考图完全一致，不得改变人物外貌；角色表情设定图，同一角色6个表情水平并排一行，1×6横向排列，半身像，正面平视，纯色浅灰背景，写实风格，电影感光影，高清细节，光线机位构图一致，仅表情变化，无水印无文字：
	\end{quote}
	
	\subsection*{统一后缀（固定加在每条结尾）}
	\begin{quote}
		\small
		，严格保持与人物图一致的外貌与服装，6个表情等宽等距横向排列，表情夸张清晰可辨
	\end{quote}
	
	\subsection*{6 个表情固定顺序（所有角色统一）}
	\begin{enumerate}[leftmargin=2em, itemsep=0.2em]
		\item 平静/常态
		\item 微笑
		\item 愤怒
		\item 惊讶
		\item 悲伤/不安
		\item 标志性表情
	\end{enumerate}
	
	\newpage
	
	\section{校长}
	
	\begin{quote}
		\small
		严格参考所提供的人物图，人物长相、服装、发型、色彩、比例与参考图完全一致，不得改变人物外貌；角色表情设定图，同一角色6个表情水平并排一行，1×6横向排列，半身像，正面平视，纯色浅灰背景，写实风格，电影感光影，高清细节，光线机位构图一致，仅表情变化，无水印无文字：从左到右6个表情——1平静（嘴角平直，眼神审视，微微仰头）、2微笑（嘴角挂一抹可随时收放的礼貌微笑，眼神不达眼底）、3愤怒（眉头紧压，嘴角下撇，眼神凌厉，面部绷紧）、4惊讶（眉毛上扬，眼睛微睁，嘴角僵住，强作镇定）、5不安（微笑消失，眼神闪躲，额头微汗，整理领带的小动作）、6标志性表情（双手撑讲台，微微仰头，标准微笑扫视全场），严格保持与人物图一致的外貌与服装，6个表情等宽等距横向排列，表情夸张清晰可辨
	\end{quote}
	
	\section{老师（开明老师）}
	
	\begin{quote}
		\small
		严格参考所提供的人物图，人物长相、服装、发型、色彩、比例与参考图完全一致，不得改变人物外貌；角色表情设定图，同一角色6个表情水平并排一行，1×6横向排列，半身像，正面平视，纯色浅灰背景，写实风格，电影感光影，高清细节，光线机位构图一致，仅表情变化，无水印无文字：从左到右6个表情——1平静（眼神犹豫，心事重重）、2微笑（勉强的温和微笑，眼神带不安）、3愤怒（眉头紧锁，嘴唇紧抿，压抑愤懑，身体试图坐直却僵硬）、4惊讶（眼睛睁大，嘴唇微张，欲言又止，身体后缩）、5悲伤不安（慢慢低头，眼神下垂，肩膀塌下，无力）、6标志性表情（欲言又止，嘴唇微动，目光躲闪，被注视时低头），严格保持与人物图一致的外貌与服装，6个表情等宽等距横向排列，表情夸张清晰可辨
	\end{quote}
	
	\section{教育部部长}
	
	\begin{quote}
		\small
		严格参考所提供的人物图，人物长相、服装、发型、色彩、比例与参考图完全一致，不得改变人物外貌；角色表情设定图，同一角色6个表情水平并排一行，1×6横向排列，半身像，正面平视，纯色浅灰背景，写实风格，电影感光影，高清细节，光线机位构图一致，仅表情变化，无水印无文字：从左到右6个表情——1平静（双手背身后，眼神审视）、2微笑（笑容来得快，眼角堆笑，带满意审视）、3愤怒（笑容凝固缓缓转头，眉头压低，眼神冷硬）、4惊讶（眉毛高挑，眼睛睁圆，嘴巴微张，身体前倾）、5满意得意（拍肩姿态，笑容满面，微微前倾）、6标志性表情（提问时微微前倾，笑容半收，眼神在审视与满意间切换），严格保持与人物图一致的外貌与服装，6个表情等宽等距横向排列，表情夸张清晰可辨
	\end{quote}
	
	\section{秘书}
	
	\begin{quote}
		\small
		严格参考所提供的人物图，人物长相、服装、发型、色彩、比例与参考图完全一致，不得改变人物外貌；角色表情设定图，同一角色6个表情水平并排一行，1×6横向排列，半身像，正面平视，纯色浅灰背景，写实风格，电影感光影，高清细节，光线机位构图一致，仅表情变化，无水印无文字：从左到右6个表情——1平静（表情中性，站姿恭敬克制）、2微笑（极浅职业性微笑，克制顺从）、3紧张（眉头微蹙，眼神紧张，双手握相机，身体前倾待命）、4惊讶（眼睛微睁，嘴巴微张，随即恢复中性）、5专注（紧盯取景框，表情严肃，手指按快门）、6标志性表情（手持相机随时拍摄，眼神跟随前方，动作精准），严格保持与人物图一致的外貌与服装，6个表情等宽等距横向排列，表情夸张清晰可辨
	\end{quote}
	
	\section{学生A}
	
	\begin{quote}
		\small
		严格参考所提供的人物图，人物长相、服装、发型、色彩、比例与参考图完全一致，不得改变人物外貌；角色表情设定图，同一角色6个表情水平并排一行，1×6横向排列，半身像，正面平视，纯色浅灰背景，写实风格，电影感光影，高清细节，光线机位构图一致，仅表情变化，无水印无文字：从左到右6个表情——1平静（标准固定微笑，眼神明亮却缺乏变化）、2标准微笑（嘴角弧度标准，直视前方）、3紧张（笑容略僵，眼神闪动，嘴唇紧抿）、4惊讶（眼睛睁大，嘴巴微张，随即恢复标准微笑）、5兴奋抢答（抢先跨步，眼神发亮，表情用力而机械）、6标志性表情（保持标准微笑一动不动，被拍肩时笑容不变），严格保持与人物图一致的外貌与服装，6个表情等宽等距横向排列，表情夸张清晰可辨
	\end{quote}
	
	\section{群众领导、群众老师（群体角色）}
	
	\begin{quote}
		\small
		严格参考所提供的人物图，群体人物的长相、服装、发型、色彩、比例与参考图完全一致，不得改变人物外貌；角色表情设定图，同一群体角色6个表情水平并排一行，1×6横向排列，多人整齐排列半身像，正面平视，纯色浅灰背景，写实风格，电影感光影，高清细节，光线机位构图一致，仅表情变化，无水印无文字：从左到右6个表情——1平静（面部严肃空洞，视线统一朝前）、2同步微笑（整齐划一笑容同时出现，眼神空洞）、3严肃（眉头微压，坐姿笔直，身体同步倾斜四十五度）、4惊讶（全体同步睁眼微张嘴，头部整齐转向同一方向）、5鼓掌（整齐拍手三下，节奏精准，表情机械）、6标志性表情（身体统一向一侧倾斜四十五度，整齐转头，同步笑声哈哈哈，表情空洞统一），严格保持与人物图一致的外貌与服装，6个表情等宽等距横向排列，表情夸张清晰可辨
	\end{quote}
	
	\section{群众学生（群体角色）}
	
	\begin{quote}
		\small
		严格参考所提供的人物图，群体人物的长相、服装、发型、色彩、比例与参考图完全一致，不得改变人物外貌；角色表情设定图，同一群体角色6个表情水平并排一行，1×6横向排列，多名青少年整齐排列半身像，正面平视，纯色浅灰背景，写实风格，电影感光影，高清细节，光线机位构图一致，仅表情变化，无水印无文字：从左到右6个表情——1平静（双手整齐放桌面，面无表情，视线统一朝讲台）、2同步微笑（全体嘴角同时上扬，弧度一致，眼神空洞）、3严肃（面无表情，坐姿笔直统一）、4惊讶（听到滴声同时抬头，眼睛同时睁大，动作如机器人一致）、5起立（同时起立，动作精准一致，表情空白）、6标志性表情（听到滴声同时低头通电，随后同时起立迈机器人步伐，面无表情，动作高度统一），严格保持与人物图一致的外貌与服装，6个表情等宽等距横向排列，表情夸张清晰可辨
	\end{quote}
	
\end{document}